\documentclass[11pt,a4paper]{article}

% ---------- Packages ----------
\usepackage[utf8]{inputenc}
\usepackage[T1]{fontenc}
\usepackage[margin=2.5cm]{geometry}
\usepackage{amsmath,amssymb,amsthm}
\usepackage{graphicx}
\usepackage{booktabs}
\usepackage[numbers,sort&compress]{natbib}
\usepackage[hidelinks]{hyperref}

\graphicspath{{figures/}}
\newtheorem{theorem}{Theorem}

% ---------- Metadata ----------
\title{A Descriptive Title for Your Paper}
\author{
  Your Name\thanks{Corresponding author: \texttt{you@example.com}} \\
  Department, University \\
}
\date{\today}

\begin{document}
\maketitle

\begin{abstract}
Summarise the problem, your approach, the key results and why they matter
in roughly 150--250 words.
\end{abstract}

\noindent\textbf{Keywords:} keyword one, keyword two, keyword three

\section{Introduction}
\label{sec:intro}

Motivate the problem and summarise related work~\cite{knuth1984,lamport1994}.
Our contributions are:
\begin{itemize}
  \item the first contribution;
  \item the second contribution.
\end{itemize}
The rest of the paper is organised as follows: Section~\ref{sec:method}
describes the method, Section~\ref{sec:results} the results and
Section~\ref{sec:conclusion} concludes.

\section{Method}
\label{sec:method}

Equations are numbered automatically:
\begin{equation}
  \mathcal{L}(\theta) = \frac{1}{N}\sum_{i=1}^{N} \ell\bigl(f_\theta(x_i), y_i\bigr).
  \label{eq:loss}
\end{equation}

\begin{theorem}
If Equation~\eqref{eq:loss} is convex in $\theta$, every local minimum is global.
\end{theorem}


\section{Results}
\label{sec:results}

Table~\ref{tab:results} compares our approach with a baseline, and
Figure~\ref{fig:overview} gives an overview.

\begin{table}[ht]
  \centering
  \caption{Example results table.}
  \label{tab:results}
  \begin{tabular}{lcc}
    \toprule
    Method   & Accuracy (\%) & Time (s) \\
    \midrule
    Baseline & 81.2 & 12.4 \\
    Ours     & \textbf{89.7} & \textbf{9.8} \\
    \bottomrule
  \end{tabular}
\end{table}

\begin{figure}[ht]
  \centering
  \includegraphics[width=0.6\linewidth]{placeholder.png}
  \caption{Replace \texttt{figures/placeholder.png} with your own figure.}
  \label{fig:overview}
\end{figure}

\section{Conclusion}
\label{sec:conclusion}

State the main contributions and future work.

\bibliographystyle{plainnat}
\bibliography{references}

\end{document}
